\documentclass[11pt]{article}
\usepackage[margin=0.85in]{geometry}
\usepackage{graphicx}
\usepackage{amsmath}
\title{R101: the marked-23 source-order fiber at $x=1$}
\author{Koblitz modular IC research}
\date{5 November 2026}
\begin{document}
\maketitle
The SHA-256-bound R51 triangular source specializes to
\[
 f_{51}(1,y)=y^3(y+1)^4,\qquad
 G_{51}(1,y,l)=l^{16}+l^8+1=(l^2+l+1)^8.
\]
Every one of the 28 rational lambda coefficients was accounted for.
The lambda derivative vanishes identically at $x=1$. Singular checked
both base-point lambda factorizations, full products, and changed-product
controls. A fresh process rebound the original donor and reproduced all
exports byte-for-byte. Independent GF(2) arithmetic and Lean 4.34.1
checked the displayed identities and length arithmetic.

\begin{center}
\includegraphics[width=0.88\linewidth]{diagram.pdf}
\end{center}

For the two-variable monogenic order $\mathbf F_2[x,y,l]/(f_{51},G_{51})$,
the $x=1$ fiber has two maximal ideals with $\mathbf F_4$ residue field.
Their local $\mathbf F_2$ dimensions are 48 and 64, totaling 112. These
are source-order lengths. The linear $N_{51}$ coordinate was not used;
local normalization is required for branch counts and valuations.

The first two preparations were retained as failures: a constant-term
parser error and an incorrect expectation that the derivative would stay
nonzero. The compatible independent verifier passed after its Python 3.9
bit-count fix. The raw archive SHA-256 is
\texttt{6d672fe4f5ea9dc4ff66d7334586959cb120bb09e7521332bfcc85e2}\linebreak
\texttt{cc731c99}.
\end{document}
